\chapter{ISO 27001 Control Mapping}
\chaptermark{ISO 27001 Mapping}
\label{ch:iso27001}

This chapter maps the secure online banking platform to ISO/IEC~27001 and summarizes gaps and compensating controls.

\section{Methodology}

To align the banking platform with widely used security practice, we applied a three-step mapping process:

\begin{description}
  \item[Inventory of existing controls:]
    Identify defense-in-depth mechanisms already implemented in the Azure architecture.
  \item[Annex A alignment:]
    Map those mechanisms to the 93 controls in ISO/IEC~27001:2022 Annex~A.
  \item[Gap analysis and compensation:]
    Identify cases where controls might be bypassed (for example zero-day exploits or token theft) and define compensating safeguards.
\end{description}

\section{Control mapping and gap analysis}

The complete mapping of technical controls to ISO/IEC~27001:2022 Annex~A---including gap analysis and compensating controls---is provided in the PDF insert on the following pages. This insert is the authoritative record for this section.

\begin{figure}[htbp]
  \centering
  \includegraphics[width=\textwidth,height=0.82\textheight,keepaspectratio,page=1]{ISO_Control_Mapping.pdf}
  \caption{ISO control mapping (full document, page 1 of 2).}
  \label{fig:iso-map-1}
\end{figure}

\begin{figure}[htbp]
  \centering
  \includegraphics[width=\textwidth,height=0.82\textheight,keepaspectratio,page=2]{ISO_Control_Mapping.pdf}
  \caption{ISO control mapping (full document, page 2 of 2).}
  \label{fig:iso-map-2}
\end{figure}

\section{Strategic outcomes}

Mapping the architecture to these controls establishes a strong baseline for digital banking. Compensating measures such as HSTS and hash-protected logs reduce residual risk and keep the system resilient when primary controls face advanced threats.
